\section{Geometry of a signed boundary representative}
\label{proj:sg:section}

The next statement isolates the extension argument needed in the
induction.  Its boundary is reduced, but its prescribed representative
need not be effective.

\begin{theorem}[Signed representative]\label{proj:thm:signed}
Let $(X,D)$ be a projective $\Q$-factorial dlt pair over $\C$, with
$D=\sum_iD_i$ reduced.  Suppose that
\[
 L=K_X+D\quad\text{is nef},\qquad L-cD\quad\text{is pseudo-effective}
 \quad\text{for some }c>0,
\]
and that
\[
 L\sim_{\Q}G=\sum_i a_iD_i,\qquad a_i\in\Q.
\]
If $L|_D$ is semiample as a rational line bundle on the reduced scheme
$D$, then $L$ is abundant:
\[
 \kappa(X,L)=\nu(X,L).
\]
\end{theorem}

The proof occupies this section, the Hodge-theoretic lifting argument
of Proposition~\ref{proj:prop:formal-lifting}, and the descent argument in
Appendix~\ref{proj:sg:descent}.  No Iitaka subadditivity assumption is used
in this theorem.

\subsection{The positive boundary and its small intersections}

Write
\[
 G=G_+-G_-,\qquad D_0=\sum_{a_i=0}D_i,
\]
where $G_+$ and $G_-$ are effective with disjoint prime supports.
Choose a positive integer $m$ such that $mG_\pm$, $mD_0$, and $mL$
are integral Cartier divisors and
\[
 N_0=\OO_X(mG)\simeq\OO_X(mL).
\]
Increase $m$ so that $N_0|_D$ is generated by global sections.  Fix the
resulting morphism
\[
 \phi:D\longrightarrow P=\PP^s,\qquad
 N_0|_D\simeq\phi^*\OO_P(1),
\]
and let $s_0$ be the rational section of $N_0$ with divisor $mG$.
Set $n=\dim X$ and $v=\nu(X,L)$, and fix an ample Cartier divisor $H$.

\begin{lemma}\label{proj:sg:numerical-boundary}
If $v=n$, then $L$ is big.  If $v=0$, then $D=0$ and
$L\sim_{\Q}0$.  In the remaining case $0<v<n$, put $r=v-1$.
Every component of $D$ has $\phi$-image of dimension at most $r$,
and some component of $G_+$ has image of dimension $r$.
Moreover,
\[
 \dim\phi\bigl(\Supp G_+\cap(\Supp G_-\cup D_0)\bigr)<r;
\]
when $r=0$, the intersection in this formula is empty.
\end{lemma}

\begin{proof}
The assertion for $v=n$ is the numerical criterion for bigness of a
nef divisor.  Suppose $v<n$.  Intersecting the pseudo-effective class
$L-cD$ with a product of nef classes gives
\[
 0\le (L-cD)L^vH^{n-v-1}
    =-c\sum_iD_iL^vH^{n-v-1}.
\]
Every summand on the right is nonnegative.  Thus each is zero.
For $v=0$, ampleness gives $D=0$, and the signed representative gives
$L\sim_{\Q}0$.

Now suppose $v>0$.  On a component of $D$, the numerical dimension
of $L|_D$ equals its image dimension under $\phi$.  The preceding
vanishing gives the upper bound $r$.  On the other hand,
\[
 0<L^vH^{n-v}=\sum_i a_iD_iL^rH^{n-v},
\]
so a positive-coefficient component attains that bound.

Consider the symmetric matrix
\[
 Q_{ij}=D_iD_jL^rH^{n-r-2}.
\]
Its off-diagonal entries are nonnegative: distinct effective
$\Q$-Cartier divisors have effective intersection cycles.  The mixed
Hodge index theorem, obtained by approximation with ample classes
and the ordinary Hodge index theorem on complete-intersection
surfaces, says that the ambient intersection form has at most one
positive direction.  In this form, $L$ is isotropic, is orthogonal
to every $D_i$, and has strictly positive pairing with $H$.
Its orthogonal space therefore has negative semidefinite form.
The same applies after pullback to a resolution, so this argument
does not require $X$ to be smooth.

The signed relation implies $Q(a_i)=0$.  Write $a=a^+-a^-$
temporarily for the positive and negative coefficient vectors.
Then
\[
 Q(a^+,a^+)=Q(a^+,a^-)\ge0.
\]
Negative semidefiniteness makes both sides zero and gives
$Qa^+=0$.  The vanishing rows belonging to negative or zero
coefficients show that each intersection of a positive component
with a negative or zero component has zero $L^r$-degree against
$H^{n-r-2}$.  Semiampleness on $D$ gives image dimension less than
$r$.  If $r=0$, an effective nonzero intersection has positive
ample degree, proving emptiness.
\end{proof}

We henceforth work in the case $0<v<n$, and set $N=n-1$.

\subsection{A geometric package for infinitesimal lifting}

The root gerbe
\[
 \mathcal P=\sqrt[\ell]{\OO_P(1)}
\]
parametrizes $\ell$-th roots of the indicated line bundle, without
a chosen section.  Its tautological line is denoted by $C$, so
$C^\ell$ is the pullback of $\OO_P(1)$.  We use the same notation
for further pullbacks of $C$.

\begin{proposition}\label{proj:prop:signed-construction}
Choose $\ell$ divisible by $m$ and every nonzero integer $|ma_i|$,
and set $a=\ell/m$.  There is a normal tame Deligne--Mumford stack
$\mathcal X$, considered on a neighborhood of a reduced Cartier
divisor $S$ of pure dimension $N=n-1$, together with a morphism
to $X$ and a reduced boundary
$T$ having no component in common with $S$, with the following
properties.
\begin{enumerate}[label=\textup{(\roman*)}]
\item
The pair $(\mathcal X,S+T)$ is log canonical, and a fixed
isomorphism of reflexive sheaves is given by
\[
 \omega_{\mathcal X}(S+T)\simeq\OO_{\mathcal X}(aS).
 \tag{\(\ast_{\mathrm{adj}}\)}\label{proj:sg:ambient-adjunction}
\]
The support of $T$ is locally set-theoretically principal.
The ambient space and $S$ are Cohen--Macaulay on scheme charts.
Off $T$, the ambient space is Gorenstein and the displayed
isomorphism is ordinary Cartier adjunction data.

\item
Put $J=(S\cap T)_{\mathrm{red}}$.  Both $S$ and $J$ are Du Bois
on scheme charts, the ideal $\mathcal I_{J/S}$ is maximal
Cohen--Macaulay, and
\[
 \mathcal B:=
 \mathcal Hom_S(\mathcal I_{J/S},\omega_S)
 \simeq\OO_S(aS).
\]
The identification agrees off $J$ with the residue convention
specified by \eqref{proj:sg:ambient-adjunction}.

\item
There is a finite representable morphism
\[
 S\longrightarrow D\times_P\mathcal P
\]
whose image covers $\Supp G_+$.  The line $\OO_S(S)$ is the pullback
of $C$.  The image of $J$ in $P$ has dimension less than $r$,
whereas the image of $S$ has dimension $r$.

\item
There is a smooth projective bundle $p:\mathcal Y\to\mathcal P$
and a representable projective morphism $g:S\to\mathcal Y$.
Writing $h=p\circ g$ and
$\mathcal T=\Spec_{\mathcal P}h_*\OO_S$, the morphism $g$ factors
through a closed embedding $\mathcal T\hookrightarrow\mathcal Y$,
and
\[
 g_*\OO_S=\OO_{\mathcal T},\qquad
 \OO_S(S)=g^*C,\qquad \mathcal B=g^*C^a.
\]

\item
For a separated quasi-projective scheme chart
$U\to\mathcal Y$ that is etale and of finite type, put
$S_U=S\times_{\mathcal Y}U$.
The etale morphism $S_U\to S$ extends compatibly to every
nilpotent thickening $qS$ in $\mathcal X$.
The resulting $qS_U$ are quasi-projective schemes, separated
and quasi-finite over $X$.
\end{enumerate}
All the adjunction and duality identifications are compatible
with changes of etale chart.
\end{proposition}

\begin{proof}
Start with the normalized simultaneous root stack of the Cartier
divisors $mG_+$, $mG_-$, and $mD_0$, taking an $\ell$-th root of
each divisor with its section.  Empty divisors cause no difficulty.
Write $S_+,S_-,S_0$ for the tautological root divisors.
To check their reducedness, work at a generic prime where the
downstairs multiplicity is $b$.  A normalized chart of
$y^\ell=x^b$ has ramification index $\ell/b$, and $y$ has order
one.  This applies to $D_0$ because $m$ divides $\ell$.
The divisors are Cartier; normality and their generic
reducedness imply that they are reduced.

Log ramification gives a log canonical pair with this full
reduced boundary and
\[
 K+S_++S_-+S_0\sim_{\Q}a(S_+-S_-).
\]
Its ambient charts are klt.  Indeed they are klt off the boundary,
and decreasing all boundary coefficients slightly removes all
zero-discrepancy places, as can be checked on a log resolution.
The integral Weil class given by the difference of the two
sides is torsion.  Take the finite representable cover formed
from its reflexive powers, with a chosen periodicity
isomorphism, and normalize.  On codimension-one charts this is
the usual cover of a torsion line bundle, hence is etale there.
The adjoint relation becomes a fixed linear equivalence;
log canonicity, klt ambient singularities, and reduced Cartier
boundary divisors persist.
More explicitly, before taking this cover write
$B=S_++S_-+S_0$.  The torsion reflexive sheaf is
\[
 \mathcal F=
 \bigl(\omega(B)\otimes\OO(-a(S_+-S_-))\bigr)^{**}.
\]
The relative spectrum of its reflexive-power algebra, with a
chosen periodicity $\mathcal F^{[q]}\simeq\OO$, has a
tautological evaluation trivializing the pulled-back
$\mathcal F$ in codimension one.  This is a morphism of
sheaves on the cover stack itself.  The resulting adjoint
isomorphism is therefore equivariant on every atlas, rather
than a separately chosen nonequivariant trivialization.

Blow up the ideal of $S_+\cap S_-$ and normalize.  Locally the
ideal has two generators, so the ordinary blowup embeds in a
relative projective line.  Its fibers, and those of its finite
normalization, have dimension at most one.  Every exceptional
divisor therefore lies over a codimension-two component of the
intersection.  Such a component is a dlt stratum downstairs and
is generically SNC.  Separate normalized Kummer charts, followed
by purity for the index cover at the smooth generic locus,
give the same description upstairs.

It follows that the tautological exceptional divisor $E$ is
reduced Cartier and
\[
 S'_+=p_1^*S_+-E,\qquad S'_-=p_1^*S_--E
\]
are disjoint reduced Cartier divisors, where $p_1$ denotes this
normalized blowup.  No codimension-two two-branch stratum is
contained in $S_0$, so $p_1^*S_0$ is reduced Cartier without an
exceptional component.  The boundary
\[
 S'_++S'_-+E+p_1^*S_0
\]
is crepant and log canonical.  Write $B'$ for this total
boundary on a chart $V'$.  The crepant equality establishes
log canonicity for every valuation.  Outside $B'$, the
blowup is an isomorphism to the klt complement of the old
boundary.  Thus a divisor $F$ with
$a(F;V',B')=0$ has center in $\Supp B'$.  As $B'$ is
effective Cartier, $\ord_F(B')>0$, and
\[
 a(F;V',0)=a(F;V',B')+\ord_F(B')>0.
\]
Divisors with positive pair discrepancy remain positive
after dropping the boundary.  This proves that $V'$ is
klt, including at higher-codimension singular loci.  Set
$S=S'_+$ and $T=S'_-+E+p_1^*S_0$ for the moment.
Near $S$, the disjoint divisor $S'_-$ contributes its unit
section, and the fixed linear equivalence becomes
\[
 \omega(S+T)\simeq\OO(aS).
\]
The ambient charts remain klt and hence Cohen--Macaulay.
Since all displayed boundary divisors are Cartier, this
isomorphism also makes those charts Gorenstein.

The strict divisor $S$ is finite over $S_+$.
Before normalization it lies in the zero section of the
projective ratio coordinate: its fibers over $S_+$ have at
most one point.  The morphism is proper, and normalization is
finite.  In particular a codimension-one point of $S\cap T$
lies over a codimension-two intersection downstairs.  The
generic SNC description above shows that $T|_S$ is generically
reduced.  It is Cartier on the Cohen--Macaulay scheme $S$,
so it is reduced everywhere.

We next remove the root characters that will not be used.
The line $\OO(S-S'_-)$ has $\ell$-th power equal to the pullback
of $N_0$.  It defines a morphism to the gerbe of $\ell$-th roots
of $N_0$ on $X$.  Take the relative coarse space over this
gerbe and call it $\mathcal X$.  On a trivializing chart this
means quotienting by the kernel of the finite group character
acting on the indicated root line.  These are ordinary
quasi-projective quotient schemes: the preceding covers are
finite, the blowup is projective, and a finite quotient preserves
quasi-projectivity.  The chartwise construction patches.

Retain $S,T$ for their reduced images.  Near $S$, the pole
section is a unit, so both the line of $S$ and its section
descend through this kernel quotient.  Thus $S$ remains
Cartier.  A finite-group norm of a local equation for the
upstairs $T$ shows that its image is set-theoretically
principal; we do not need reduced $T$ to be Cartier.
The only divisorial ramification is on the boundary.
Log ramification therefore preserves the lc pair and descends
the equivariant adjoint isomorphism to
\eqref{proj:sg:ambient-adjunction}.  Cohen--Macaulayness of the
ambient chart and of $S$ follows from the finite CM covers
and their invariant direct summands.  For clarity, the
coarse kernel acts trivially on the retained root line,
and hence on $\OO(aS)$.  Equivariance of the fixed adjoint
isomorphism gives the same kernel character on
$\omega(S+T)$.  Taking invariants therefore descends the
isomorphism, with log ramification identifying the
invariant log dualizing sheaf.  Off $T$ the descended
adjoint isomorphism makes the canonical sheaf invertible,
giving the claimed Gorenstein property there.

The reduced supports $S$ and $J$ are unions of log canonical
centers: intersections of lc centers are again unions of
lc centers.  Hence both are Du Bois by
\cite[Theorems~1.4 and~1.7]{Proj-KollarKovacs2010}.
For the more precise coherent statement, on one quotient
chart write $S^{\mathrm{up}}\to S$ for the finite cover used
above.  Since $T^{\mathrm{up}}|_{S^{\mathrm{up}}}$ is reduced,
\[
 \mathcal I_{J/S}
 =
 \bigl(f_*\OO_{S^{\mathrm{up}}}
       (-T^{\mathrm{up}}|_{S^{\mathrm{up}}})\bigr)^{\mathrm{inv}}.
\]
Indeed an invariant function vanishes on the reduced quotient
image precisely when its pullback vanishes on this reduced
preimage.  The sheaf on the right is maximal Cohen--Macaulay:
upstairs it is invertible on a CM scheme, finite pushforward
has the same depth over the quotient, and invariants are a
direct summand in characteristic zero.

Finite-map duality with trace, followed by invariants, gives
\[
 \mathcal Hom_S(\mathcal I_{J/S},\omega_S)
 =
 \bigl(f_*\omega_{S^{\mathrm{up}}}
       (T^{\mathrm{up}}|_{S^{\mathrm{up}}})\bigr)^{\mathrm{inv}}
 \simeq\OO_S(aS).
\]
This use of duality allows ramification and does not assert
that $\omega_S$ itself is invertible.  The last isomorphism
is upstairs Cartier adjunction followed by descent of the
$S$ line.  Off $J$ it is the ordinary residue determined by
the fixed ambient isomorphism.  Normalized trace preserves
that convention, proving compatibility on overlaps.

Restricted to $S$, the root gerbe of $N_0$ is
$D\times_P\mathcal P$.  The strict-divisor finiteness already
proved, together with removal of the relative inertia kernel,
makes $S\to D\times_P\mathcal P$ finite and representable.
Its image covers the positive boundary, and its root line is
$\OO_S(S)=C$.  Lemma~\ref{proj:sg:numerical-boundary} gives all the
asserted image-dimension bounds, including the bound for $J$.
The resulting $h:S\to\mathcal P$ is representable projective.

Its Stein algebra defines a finite stack
$\mathcal T=\Spec_{\mathcal P}h_*\OO_S$ over $\mathcal P$.
There are enough vector bundles on $\mathcal P$ to generate
this coherent algebra as a module.  Namely, decompose by the
finitely many inertia characters, twist each summand by a
power of $C$ to descend it to $P$, and apply Serre generation.
A vector-bundle surjection gives an embedding of $\mathcal T$
in a vector bundle over $\mathcal P$.  Its closure in the
projective completion is still $\mathcal T$, since it is
proper over the base.  This gives $\mathcal Y$ and $g$.
The Stein construction supplies $g_*\OO_S=\OO_{\mathcal T}$
and the required line identities.

Finally, etale morphisms extend uniquely over nilpotent
thickenings, compatibly as $q$ varies.  The reduction $S_U$
is a scheme.  The extension $qS_U$ has trivial inertia and
is an algebraic space.  It is separated over $X$: the chosen
chart is separated, the morphism to the root gerbe is finite,
and that gerbe has finite diagonal.  Its finite-type
geometric fibers over $X$ are zero-dimensional, and this is
unchanged by nilpotents.  Thus $qS_U\to X$ is separated and
quasi-finite.  Zariski's main theorem embeds it in a scheme
finite over $X$, proving that it is a quasi-projective
scheme.  This completes the construction.
\end{proof}
